\documentclass[11pt,a4paper]{article}

% ----- Encoding & fonts (clean sans-serif) -----
\usepackage[utf8]{inputenc}
\usepackage[T1]{fontenc}
\usepackage{helvet}
\renewcommand{\familydefault}{\sfdefault}
\usepackage{microtype}

% ----- Math + formatting -----
\usepackage{amsmath,amssymb}
\usepackage{enumitem}
\usepackage{array}
\usepackage{tabularx}
\usepackage{booktabs}

% ----- Figures + drawing -----
\usepackage{graphicx}
\usepackage{float}

% ----- Code listings -----
\usepackage{listings}
\usepackage[table]{xcolor}

% ----- Layout -----
\usepackage[margin=1in,headheight=14pt]{geometry}
\usepackage{parskip}
\usepackage{titlesec}
\usepackage{fancyhdr}
\usepackage[hidelinks]{hyperref}

% ----- Theme colours -----
\definecolor{accent}{HTML}{1F4E79}     % deep professional blue
\definecolor{accentdk}{HTML}{2E2E2E}   % near-black for subsections
\definecolor{accentlt}{HTML}{D9E2F0}   % light blue table tint
\definecolor{rulegray}{HTML}{B9B9B9}
\definecolor{codebg}{HTML}{F7F7F4}
\definecolor{codecomment}{HTML}{6A737D}
\definecolor{codekw}{HTML}{0B5FA5}
\definecolor{codestr}{HTML}{B5210F}

% ----- Section / subsection styling -----
\titleformat{\section}
  {\color{accent}\Large\bfseries}{}{0pt}{}
  [\vspace{2pt}{\color{accent}\titlerule[1.1pt]}]
\titlespacing*{\section}{0pt}{16pt}{8pt}

\titleformat{\subsection}
  {\color{accentdk}\large\bfseries}{}{0pt}{}
\titlespacing*{\subsection}{0pt}{12pt}{3pt}

% ----- Page header / footer -----
\pagestyle{fancy}
\fancyhf{}
\fancyhead[L]{\small\color{accentdk} FPGA 5 -- Keyboard}
\fancyhead[R]{\small\color{accentdk}FPGA Lab}
\fancyfoot[C]{\small\thepage}
\renewcommand{\headrulewidth}{0.4pt}
\renewcommand{\footrulewidth}{0pt}

% ----- Verilog listing style -----
\lstdefinestyle{vlog}{
    language=Verilog,
    backgroundcolor=\color{codebg},
    basicstyle=\ttfamily\footnotesize,
    keywordstyle=\color{codekw}\bfseries,
    commentstyle=\color{codecomment}\itshape,
    stringstyle=\color{codestr},
    numbers=left,
    numberstyle=\tiny\color{gray},
    numbersep=8pt,
    frame=single,
    rulecolor=\color{rulegray},
    breaklines=true,
    showstringspaces=false,
    tabsize=4,
    xleftmargin=2.2em,
    framexleftmargin=2.2em,
    aboveskip=10pt,
    belowskip=6pt,
    captionpos=b,
}
\lstset{style=vlog}

% ----- A placeholder box for screenshots still to be pasted in -----
% Replace each \placeholder{...} with a real \includegraphics{...} once the
% image is exported from Vivado / taken with a camera.
\newcommand{\placeholder}[2][0.92\linewidth]{%
  \begin{center}
  \setlength{\fboxsep}{0pt}%
  \fcolorbox{rulegray}{white}{\begin{minipage}[c][4.4cm][c]{#1}
    \centering\itshape\color{gray}
    [ paste screenshot here ]\\[5pt]
    \footnotesize #2
  \end{minipage}}
  \end{center}%
}

% ----- A real figure: framed image + italic caption underneath -----
\newcommand{\reportfig}[3][\linewidth]{%
  \begin{center}
  \setlength{\fboxsep}{1.5pt}%
  \fcolorbox{rulegray}{white}{\includegraphics[width=#1]{#2}}\\[4pt]
  {\footnotesize\itshape #3}
  \end{center}%
}

\setlength{\parindent}{0pt}
\setlength{\parskip}{0.6em}

\begin{document}

% ======================================================================
%  TITLE BLOCK
% ======================================================================
\thispagestyle{fancy}
\begin{center}
  {\color{accent}\rule{\linewidth}{1.4pt}}\\[10pt]
  {\fontsize{24}{28}\selectfont \textbf{\color{accent}FPGA Experiment Report}}\\[6pt]
  {\large Experiment 5 --- PS/2 Keyboard Interface}\\[8pt]
  {\color{accent}\rule{\linewidth}{1.4pt}}
\end{center}

\vspace{6pt}

\renewcommand{\arraystretch}{1.3}
\begin{center}
\begin{tabular}{@{}l l l@{}}
\textbf{Task:} & FPGA 5 --- Keyboard (PS/2 interface) & \\
\textbf{Student 1:} & Mahmood Stitia & \textbf{ID:} 214032682 \\
\textbf{Student 2:} & Saleh Khalil  & \textbf{ID:} 213310485 \\
\end{tabular}
\end{center}

\vspace{6pt}

% ======================================================================
%  PART 1 - STATUS UPDATE
% ======================================================================
\section*{Part 1 --- Status update}

The table below lists the submitted files and their status. The status is one of
\emph{Done} (module and simulation written and passing), \emph{Partially done}
(started but not fully working --- explained in the notes), or \emph{Not started}.

\vspace{4pt}
\renewcommand{\arraystretch}{1.5}
\begin{tabularx}{\textwidth}{|>{\bfseries\ttfamily}l|c|X|}
\hline
\rowcolor{accent}
\textcolor{white}{\textbf{File / Step}} & \textcolor{white}{\textbf{Status}} & \textcolor{white}{\textbf{Notes / comments}} \\ \hline
\rowcolor{accentlt}
Ps2\_Interface.v     & Done & \\ \hline
Ps2\_Interface\_tb.v & Done & \\ \hline
\rowcolor{accentlt}
Ps2\_Display.v       & Done & \\ \hline
Ps2\_Top.v           & Done & \\ \hline
\rowcolor{accentlt}
task1.xdc            & Done & \\ \hline
\end{tabularx}
\renewcommand{\arraystretch}{1.0}

\subsection*{Issues}

Any errors present at the time of submission are described here (per module /
simulation, with what the error means).

\begin{itemize}[topsep=2pt,itemsep=2pt,leftmargin=1.4em]
  \item \texttt{Ps2\_Interface.v}:
  \item \texttt{Ps2\_Interface\_tb.v}:
  \item \texttt{Ps2\_Display.v}:
  \item \texttt{Ps2\_Top.v}:
  \item \texttt{task1.xdc}:
\end{itemize}

% ======================================================================
%  PART 2 - SIMULATIONS AND ANSWERS
% ======================================================================
\newpage
\section*{Part 2 --- Simulations and answers}
\vspace{1cm}
\subsection*{Simulation --- Ps2\_Interface}








To verify the \texttt{Ps2\_Interface} module, we  emulate the behavior of a real PS/2 keyboard. The PS/2 clock is generated with a
\SI{60}{\micro\second} period (a \SI{30}{\micro\second} half-period),
which corresponds to a frequency of approximately \SI{16.7}{\kilo\hertz}.
This value sits inside the standard PS/2 clock range
(\SIrange{10}{16.7}{\kilo\hertz}), so the simulation exercises the
interface at a realistic device speed.

Each key event is transmitted by the \texttt{send\_byte} task, which
serializes a byte into the standard 11-bit PS/2 frame:
one start bit (\texttt{0}), eight data bits sent LSB-first, one odd-parity
bit, and one stop bit (\texttt{1}). For every bit the data line is updated
and the PS/2 clock is toggled low-then-high, reproducing the way a keyboard
shifts data out on its own clock. A short idle gap is inserted between
frames to separate consecutive presses.


\vspace{0.5cm}
\subsubsection*{Simulated Scenario}

The testbench plays out three scenarios that cover the key behaviours of the
interface:

\begin{enumerate}
    \item \textbf{Press `5' (\texttt{0x73}).}
    A single make-code is sent. We expect \texttt{keyPressed} to pulse and
    \texttt{scancode} to equal \texttt{0x73}.

    \item \textbf{Release `5' then press `0'.}
    First the break sequence \texttt{F0 73} is sent (the \texttt{F0} prefix
    marks a key release), and the interface correctly ignores it without
    raising \texttt{keyPressed}. Immediately afterwards the make-code
    \texttt{0x70} (`0') is sent, and we expect \texttt{keyPressed} to pulse
    with \texttt{scancode == 0x70}. This checks that a release of one key is
    discarded and a following new press is still detected.

    \item \textbf{Hold `0' (auto-repeat).}
    The same make-code \texttt{0x70} is sent a second time with no
    intervening release. Because the key is already considered pressed, the
    interface must \emph{not} raise \texttt{keyPressed} again, while
    \texttt{scancode} still reads \texttt{0x70}. This verifies that the
    keyboard's typematic auto-repeat does not generate spurious new presses.
\end{enumerate}


\vspace{1cm}
\subsection{Metastability and Clock-Domain Considerations}

The PS/2 clock and data lines come from the keyboard and are
\emph{asynchronous} to the FPGA. Sampling such a signal while it is
transitioning can drive a flip-flop into a \emph{metastable} state, where its
output settles to an undefined value and may corrupt downstream logic.

Following the approach from class, we only sample the line when it is
guaranteed stable. The PS/2 protocol holds \texttt{PS2Data} steady through the
falling edge of \texttt{PS2Clk}, so our receiver shifts the data in on
\texttt{negedge PS2Clk}, capturing each bit in its stable window and avoiding
metastable sampling.
















\newpage
\reportfig[\linewidth]{tb_waveform.png}{Full Ps2\_Interface\_tb
 with $TCL$ $Console$ message \textbf{Test Passed - Ps2\_Interface\_tb}}

\textbf{Signals} (top to bottom in the waveform):
\begin{itemize}[noitemsep,topsep=2pt,leftmargin=1.4em]
  \item \texttt{PS2Clk} - the keyboard clock (input). It toggles only while a byte is being sent and is idle-high.
  \item \texttt{rstn} - active-low \emph{asynchronous} reset (driven by the test-bench).
  \item \texttt{correct} - a test-bench-only pass/fail flag.
  \item \texttt{PS2Data} - the serial data line. Carries 11-bit frame, sampled on the falling edge of \texttt{PS2Clk}.
  \item \texttt{byte\_to\_send[7:0]} - a test-bench register holding the byte currently being bit-banged onto \texttt{PS2Data}.
  \item \texttt{scancode[7:0]} - interface output: the most recent accepted make-code.
  \item \texttt{keyPressed} - interface output: a one-cycle pulse on the first make-code of a new press.
  \item \texttt{parity\_ok} - internal interface signal : high when the frame's odd-parity bit matches the 8 data bits.
\end{itemize}

\reportfig[\linewidth]{tb_waveform_.pdf}{Annotated close-up of the same
simulation. The red notes mark the first key edge and that the first bit is sent
on the \emph{negative} clock edge.}

% ======================================================================
%  ELABORATED DESIGN
% ======================================================================
\newpage
\section*{Elaborated Design}

\subsection*{Top module}

\reportfig[\linewidth]{linter_elaborated_design.png}{RTL-elaborated schematic of
\texttt{Ps2\_Top}: \texttt{u\_interface} (\texttt{Ps2\_Interface}, PS2Clk domain)
feeds \texttt{keyPressed} and \texttt{scancode[7:0]} into \texttt{u\_display}
(\texttt{Ps2\_Display}, 100\,MHz domain); the \texttt{reset} push-button passes
through an \texttt{RTL\_INV} (\texttt{rstn\_i}) into \texttt{rstn}. Inputs
\texttt{clk, PS2Clk, PS2Data, reset} are on the left and outputs
\texttt{an[3:0], dp, led, seg[6:0]} on the right.}

\subsection*{Ps2\_Interface}

\reportfig[\linewidth]{sch_innterface.png}{RTL-elaborated schematic of the Interface}

% ======================================================================
%  IMPLEMENTATION
% ======================================================================
\newpage
\section*{Implementation}

\subsection*{Synthesis schematic}

\reportfig[\linewidth]{synthesis_schematic.png}{Post-synthesis schematic. Every
external port now passes through an \texttt{IBUF}/\texttt{OBUF}, the two clocks}

\subsection*{Project summary}

\reportfig[\linewidth]{project_summary.png}{Post-implementation project summary}

\subsection*{Timing summary}

\reportfig[\linewidth]{design_timing_report.png}{Design Timing Summary after
implementation}

\subsection*{Critical Path}

\reportfig[\linewidth]{critical_path.png}{Intra-Clock Paths report. The worst
paths (Path~1--4) run inside \texttt{u\_display} from
\texttt{led\_timer\_reg[16]/C} to \texttt{led\_timer\_reg[20..23]/CE}, with slack
$5.269$\,ns, 3 logic levels and $4.361$\,ns total delay ($1.054$\,ns logic
$+\,3.307$\,ns net) against the $10$\,ns requirement}
\newpage
\subsection*{Device}

\reportfig[1\linewidth]{synthesized_design.png}{Device (floorplan) view of the
placed-and-routed design on the Artix-7 (\texttt{XC7A35T})}

% ======================================================================
%  PROBLEMS WE FACED
% ======================================================================

\reportfig[1\linewidth]{system.jpeg}{The full system on hardware: the numeric
keypad plugged into the Basys-3 USB-host port. Key "5" is pressed and its
PS/2 scan-code 73 is shown on the two right-hand 7-segment digits.}


% ======================================================================
%  APPENDIX - VERILOG SOURCE CODE
% ======================================================================
\newpage
\section*{Appendix --- Verilog Source Code}

The full source of the submitted files is listed below. Paste the code of each
file inside its block.

\subsection*{Ps2\_Interface.v}
\begin{lstlisting}
module Ps2_Interface(
    input wire PS2Clk,
    input wire rstn,
    input wire PS2Data,
    output reg [7:0] scancode,
    output reg keyPressed
    );

    reg [21:0] shift_reg;
    reg [3:0]  bit_count;
    reg is_valid;
    wire parity_ok;
    wire [7:0] cur_byte  = shift_reg[20:13];
    wire [7:0] prev_byte = shift_reg[9:2];

    always @(negedge PS2Clk or negedge rstn) begin
        if (!rstn) begin
            bit_count <= 4'b0;
            shift_reg <= 22'b0;
            scancode <= 8'b0;
            keyPressed <= 1'b0;
            is_valid <= 1'b1;
        end else begin
            shift_reg <= {PS2Data, shift_reg[21:1]};
            bit_count <= (bit_count == 4'd10)? 4'd0 : bit_count + 4'd1;
            keyPressed <= 1'b0; // default

            if (bit_count == 4'd10) begin
                if (cur_byte == 8'hE0 || cur_byte == 8'hF0) begin
                    // skip
                end
                else if (prev_byte == 8'hF0) begin
                    is_valid <= 1'b1;
                end
                else begin
                    scancode <= cur_byte;
                    if (is_valid) begin
                        keyPressed <= parity_ok;
                        is_valid <= 1'b0;
                    end
                end
            end 
               
        end

    end

    assign parity_ok = (shift_reg[21] == ~(^cur_byte));

endmodule
\end{lstlisting}

\newpage
\subsection*{Ps2\_Interface\_tb.v}
\begin{lstlisting}
module Ps2_Interface_tb();

    localparam CLK_HALF = 30_000;   // 30 us half -> 60 us period -> ~16.7 kHz PS/2 clock

    reg PS2Clk, rstn, PS2Data, correct;
    wire [7:0] scancode;
    wire keyPressed;
    wire parity_ok = uut.parity_ok;
    reg [7:0] byte_to_send;


    // Instantiate the UUT (Unit Under Test)
    Ps2_Interface uut(
        .PS2Clk(PS2Clk),
        .rstn(rstn),
        .PS2Data(PS2Data),
        .scancode(scancode),
        .keyPressed(keyPressed)
        );

    initial begin

        if($test$plusargs("vcd")) begin
            $dumpfile("Ps2_Interface_tb.vcd");
            $dumpvars(0, Ps2_Interface_tb);
        end

        correct = 1;
        PS2Clk = 1;  
        PS2Data = 1; 
        rstn = 0;

        #(CLK_HALF);
        rstn = 1;
        #(CLK_HALF);

        // 1- press '5' (0x73)
        byte_to_send = 8'h73;
        send_byte(byte_to_send);
        correct = correct & keyPressed & (scancode == 8'h73);

        // 2- release '5' (F0 73) then press '0' (0x70)
        byte_to_send = 8'hF0;
        send_byte(byte_to_send);
        byte_to_send = 8'h73;
        send_byte(byte_to_send);
        byte_to_send = 8'h70;
        send_byte(byte_to_send);
        correct = correct & keyPressed & (scancode == byte_to_send);

        // 3- hold '0' (auto-repeat 0x70)
        send_byte(byte_to_send);
        correct = correct & ~keyPressed & (scancode == byte_to_send);


        if (correct)
            $display("Test Passed - %m");
        else
            $display("Test Failed - %m");
        $finish;
    end

    // send the pressed key in serial format - 11 bits
    task send_byte(input [7:0] b);
        integer k;
        reg [10:0] frame;
        begin
            frame = {1'b1, ~(^b), b, 1'b0};
            for (k = 0; k < 11; k = k + 1) begin
                PS2Data = frame[k];
                #CLK_HALF;
                PS2Clk = 0;
                #CLK_HALF;
                PS2Clk = 1;
            end
            #(CLK_HALF*4); // idle time gap between presses
        end
    endtask

endmodule
\end{lstlisting}

\newpage
\subsection*{Ps2\_Display.v}
\begin{lstlisting}
module Ps2_Display(
    input  wire       clk,         // 100 MHz system clock
    input  wire       rstn,        // active-low reset
    input  wire       keyPressed,  // pulse from Ps2_Interface (slow domain)
    input  wire [7:0] scancode,    // byte from Ps2_Interface (slow domain)
    output wire [6:0] seg,         // 7-segment cathodes (active low)
    output wire [3:0] an,          // 4 digit anodes (active low)
    output wire       dp,          // decimal point
    output reg        led          // strobe LED, one visible blink per press
    );
    reg [6:0] a_to_g;
    reg [3:0] an_reg;

    assign seg = a_to_g;
    assign an  = an_reg;
    assign dp  = 1'b1;
    
    wire [1:0] s;    
    reg  [3:0] digit;

    reg kp_sync1, kp_sync2, kp_prev;
    reg [7:0] valid_scancode;
    
    always @(posedge clk or negedge rstn) begin
        if (!rstn) begin
            kp_sync1 <= 1'b0;
            kp_sync2 <= 1'b0;
            kp_prev  <= 1'b0;
            valid_scancode <= 8'h00;
        end else begin
            kp_sync1 <= keyPressed;
            kp_sync2 <= kp_sync1;
            kp_prev  <= kp_sync2;

            if (kp_sync2 & ~kp_prev) begin
                valid_scancode <= scancode;
            end
        end
    end
    
    wire kp_edge = kp_sync2 & ~kp_prev;

    reg [23:0] led_timer;
    
    always @(posedge clk or negedge rstn) begin
        if (!rstn) begin
            led_timer <= 24'd0;
            led <= 1'b0;
        end else if (kp_edge) begin
            led_timer <= 24'd10_000_000; // ~0.1 seconds at 100MHz
            led <= 1'b1;
        end else if (led_timer > 0) begin
            led_timer <= led_timer - 1;
            led <= 1'b1;
        end else begin
            led <= 1'b0;
        end
    end

    reg [19:0] clkdiv = 20'b0;
    
    always @(posedge clk or negedge rstn) begin
        if (!rstn)
            clkdiv <= 0;
        else
            clkdiv <= clkdiv + 1;
    end
    
    assign s = clkdiv[19:18];

    always @(*) begin
        an_reg = 4'b1111; 
        
        case(s)
            2'b00: begin 
                digit = valid_scancode[3:0];
                an_reg[0] = 1'b0;
            end
            2'b01: begin 
                digit = valid_scancode[7:4];
                an_reg[1] = 1'b0;
            end
            default: begin 
                digit = 4'h0; 
                an_reg = 4'b1111;
            end
        endcase
    end

    always @(*) begin
        case(digit)
            //////////<---MSB-LSB<---/////
            4'h0: a_to_g = 7'b1000000; // 0
            4'h1: a_to_g = 7'b1111001; // 1
            4'h2: a_to_g = 7'b0100100; // 2
            4'h3: a_to_g = 7'b0110000; // 3
            4'h4: a_to_g = 7'b0011001; // 4
            4'h5: a_to_g = 7'b0010010; // 5
            4'h6: a_to_g = 7'b0000010; // 6
            4'h7: a_to_g = 7'b1111000; // 7
            4'h8: a_to_g = 7'b0000000; // 8
            4'h9: a_to_g = 7'b0010000; // 9
            4'hA: a_to_g = 7'b0001000; // A
            4'hB: a_to_g = 7'b0000011; // b
            4'hC: a_to_g = 7'b1000110; // C
            4'hD: a_to_g = 7'b0100001; // d
            4'hE: a_to_g = 7'b0000110; // E
            4'hF: a_to_g = 7'b0001110; // F
            default: a_to_g = 7'b1111111; // blank
        endcase
    end
endmodule
\end{lstlisting}

\newpage
\subsection*{Ps2\_Top.v}
\begin{lstlisting}
module Ps2_Top(
    input  wire       clk,        // W5  - 100 MHz system clock
    input  wire       reset,      // U18 - btnC, active high (pressed = 1)
    input  wire       PS2Clk,     // C17 - keyboard clock
    input  wire       PS2Data,    // B17 - keyboard data
    output wire [6:0] seg,        // 7-segment cathodes
    output wire [3:0] an,         // 7-segment anodes
    output wire       dp,         // decimal point
    output wire       led         // U16 - strobe LED (LD0)
    );

    wire rstn = ~reset;
    wire keyPressed;
    wire [7:0] scancode;

    Ps2_Display u_display(
    .clk  (clk),
    .rstn (rstn),
    .keyPressed (keyPressed),
    .scancode   (scancode),
    .seg  (seg),
    .an   (an),
    .dp   (dp),
    .led  (led));

    Ps2_Interface u_interface(
    .PS2Clk   (PS2Clk),
    .rstn     (rstn),
    .PS2Data  (PS2Data),
    .scancode (scancode),
    .keyPressed (keyPressed));

endmodule
\end{lstlisting}

\newpage
\subsection*{task1.xdc}
\begin{lstlisting}[language=]
# Clock signal
set_property PACKAGE_PIN W5 [get_ports clk]
set_property IOSTANDARD LVCMOS33 [get_ports clk]
create_clock -period 10.000 -name sys_clk_pin -waveform {0.000 5.000} -add [get_ports clk]
# LEDs
set_property PACKAGE_PIN U16 [get_ports led]
set_property IOSTANDARD LVCMOS33 [get_ports led]
#7 segment display
set_property PACKAGE_PIN W7 [get_ports {seg[0]}]
set_property IOSTANDARD LVCMOS33 [get_ports {seg[0]}]
set_property PACKAGE_PIN W6 [get_ports {seg[1]}]
set_property IOSTANDARD LVCMOS33 [get_ports {seg[1]}]
set_property PACKAGE_PIN U8 [get_ports {seg[2]}]
set_property IOSTANDARD LVCMOS33 [get_ports {seg[2]}]
set_property PACKAGE_PIN V8 [get_ports {seg[3]}]
set_property IOSTANDARD LVCMOS33 [get_ports {seg[3]}]
set_property PACKAGE_PIN U5 [get_ports {seg[4]}]
set_property IOSTANDARD LVCMOS33 [get_ports {seg[4]}]
set_property PACKAGE_PIN V5 [get_ports {seg[5]}]
set_property IOSTANDARD LVCMOS33 [get_ports {seg[5]}]
set_property PACKAGE_PIN U7 [get_ports {seg[6]}]
set_property IOSTANDARD LVCMOS33 [get_ports {seg[6]}]

set_property PACKAGE_PIN V7 [get_ports dp]
set_property IOSTANDARD LVCMOS33 [get_ports dp]

set_property PACKAGE_PIN U2 [get_ports {an[0]}]
set_property IOSTANDARD LVCMOS33 [get_ports {an[0]}]
set_property PACKAGE_PIN U4 [get_ports {an[1]}]
set_property IOSTANDARD LVCMOS33 [get_ports {an[1]}]
set_property PACKAGE_PIN V4 [get_ports {an[2]}]
set_property IOSTANDARD LVCMOS33 [get_ports {an[2]}]
set_property PACKAGE_PIN W4 [get_ports {an[3]}]
set_property IOSTANDARD LVCMOS33 [get_ports {an[3]}]


#Buttons
set_property PACKAGE_PIN U18 [get_ports reset]
set_property IOSTANDARD LVCMOS33 [get_ports reset]
#USB HID (PS/2)
set_property PACKAGE_PIN C17 [get_ports PS2Clk]
set_property IOSTANDARD LVCMOS33 [get_ports PS2Clk]
set_property PULLUP true [get_ports PS2Clk]
set_property PACKAGE_PIN B17 [get_ports PS2Data]
set_property IOSTANDARD LVCMOS33 [get_ports PS2Data]
set_property PULLUP true [get_ports PS2Data]

# Don't even think to change the following values !!!!!
set_property BITSTREAM.GENERAL.COMPRESS TRUE [current_design]
set_property BITSTREAM.CONFIG.CONFIGRATE 33 [current_design]
set_property CONFIG_MODE SPIx4 [current_design]
set_property CFGBVS VCCO [current_design]
set_property CONFIG_VOLTAGE 3.3 [current_design]
\end{lstlisting}

\end{document}
